\documentclass[11pt]{article}

\usepackage[margin=1in]{geometry}
\usepackage{amsmath, amssymb}
\usepackage{bm}
\usepackage{xcolor}
\usepackage{hyperref}
\usepackage{mdframed}
\usepackage{booktabs}

\hypersetup{
    colorlinks=true,
    linkcolor=blue!60!black,
    urlcolor=blue!60!black,
}

\newcommand{\vect}[1]{\bm{#1}}
\newcommand{\unit}[1]{\hat{\bm{#1}}}
\newcommand{\norm}[1]{\left\|#1\right\|}
\newcommand{\atantwo}{\operatorname{atan2}}

\newmdenv[
  linecolor=blue!50!black,
  linewidth=0.5pt,
  backgroundcolor=blue!3,
  roundcorner=4pt,
  skipabove=8pt,
  skipbelow=8pt,
]{keybox}

\newmdenv[
  linecolor=orange!70!black,
  linewidth=0.5pt,
  backgroundcolor=orange!5,
  roundcorner=4pt,
  skipabove=8pt,
  skipbelow=8pt,
]{warnbox}

\title{Computing Classical DH Parameters from a URDF}
\author{}
\date{}

\begin{document}
\maketitle

\section{Overview}

This document describes how to compute the classical (Spong) Denavit--Hartenberg
parameters $(a_i, \alpha_i, d_i, \theta_i)$ for each joint $i$ of a serial
manipulator, starting from a URDF file.

Each joint's DH row describes the homogeneous transform from frame $i{-}1$ to
frame $i$:
\[
  T_i^{i-1} = R_z(\theta_i)\,T_z(d_i)\,T_x(a_i)\,R_x(\alpha_i).
\]

The four parameters are:
\begin{itemize}
  \item $a_i$: link length --- distance between consecutive joint axes along the
    common normal.
  \item $\alpha_i$: link twist --- signed angle between consecutive joint
    axes, measured about $\unit{x}_i$.
  \item $d_i$: joint offset --- signed distance along $\unit{z}_{i-1}$ from the
    previous frame origin to frame $i$'s origin.
  \item $\theta_i$: joint angle offset --- signed rotation about
    $\unit{z}_{i-1}$ that aligns $\unit{x}_{i-1}$ with $\unit{x}_i$. For
    revolute joints, the runtime angle is $\theta_i + q_i$.
\end{itemize}

\begin{warnbox}
\textbf{Notation convention.} To avoid a common source of confusion, we use
two distinct symbols throughout this document:
\begin{itemize}
  \item $\unit{a}_i$ denotes the direction of joint $i$'s axis as extracted
        from the URDF (in the base frame).
  \item $\unit{z}_i$ denotes the $z$-axis of DH frame $i$.
\end{itemize}
In classical DH, frame $i$ has its $z$-axis along \emph{joint $i{+}1$'s} axis,
so $\unit{z}_i = \unit{a}_{i+1}$. To compute joint $i$'s DH row, we need the
geometry between axes $i$ and $i{+}1$ --- this is the single most common
source of derivation errors.
\end{warnbox}

\section{Step 1: Extract Axis Data from the URDF}

To compute the DH row for joint $i$, we need four values, all expressed in
the base frame: $\vect{p}_i$, $\unit{a}_i$, $\vect{p}_{i+1}$, $\unit{a}_{i+1}$.

As we walk the kinematic chain, we maintain two running quantities:
\begin{itemize}
  \item $R_{\mathrm{cum}}$ --- cumulative rotation from base to the current
        link's frame.
  \item $\vect{p}_{\mathrm{cum}}$ --- cumulative position of the current
        link's origin, in the base frame.
\end{itemize}
Initial values (at the base): $R_{\mathrm{cum}} = I$,
$\vect{p}_{\mathrm{cum}} = (0, 0, 0)$.

Suppose the URDF gives us, for each joint $i$:
\begin{itemize}
  \item \texttt{rpy} $= (r_i, \beta_i, \gamma_i)$ --- roll, pitch, yaw angles
        (URDF uses fixed-axis / extrinsic convention:
        $R = R_z(\gamma)\,R_y(\beta)\,R_x(r)$).
  \item \texttt{xyz} $= \vect{t}_i$ --- translation of the joint's origin in
        the parent link's frame.
  \item \texttt{axis} $= \unit{a}_i^{\mathrm{local}}$ --- unit direction of
        the joint's rotation axis, expressed in the child link's frame.
\end{itemize}

\paragraph{Computing $\vect{p}_1$ and $\unit{a}_1$.} Process joint 1 in
three substeps, in order:

\textbf{(a)} Update the position: rotate the joint's translation by the
\emph{current} $R_{\mathrm{cum}}$ (i.e., the parent's rotation), then add:
\[
  \vect{p}_1 = \vect{p}_{\mathrm{cum}} + R_{\mathrm{cum}} \cdot \vect{t}_1
             = (0, 0, 0) + I \cdot \vect{t}_1 = \vect{t}_1.
\]

\textbf{(b)} Update the cumulative rotation:
\[
  R_{\mathrm{cum}} \leftarrow R_{\mathrm{cum}} \cdot R(r_1, \beta_1, \gamma_1)
                            = I \cdot R(r_1, \beta_1, \gamma_1)
                            = R(r_1, \beta_1, \gamma_1).
\]

\textbf{(c)} Compute the joint axis direction, using the \emph{updated}
$R_{\mathrm{cum}}$:
\[
  \unit{a}_1 = R_{\mathrm{cum}} \cdot \unit{a}_1^{\mathrm{local}}.
\]

\paragraph{Computing $\vect{p}_2$ and $\unit{a}_2$.} After processing joint 1,
the running state is:
\[
  R_{\mathrm{cum}} = R(r_1, \beta_1, \gamma_1), \qquad
  \vect{p}_{\mathrm{cum}} = \vect{t}_1.
\]
Apply the same three substeps to joint 2.

\textbf{(a)} Update the position: rotate joint 2's translation by the
\emph{current} $R_{\mathrm{cum}}$ (i.e., joint 1's rotation, since joint 1's
child link is the parent of joint 2), then add:
\[
  \vect{p}_2 = \vect{p}_{\mathrm{cum}} + R_{\mathrm{cum}} \cdot \vect{t}_2
             = \vect{t}_1 + R(r_1, \beta_1, \gamma_1) \cdot \vect{t}_2.
\]
Unlike joint 1, this rotation is not guaranteed to be the identity, so the
components of $\vect{t}_2$ may mix into $\vect{p}_2$.

\textbf{(b)} Update the cumulative rotation:
\[
  R_{\mathrm{cum}} \leftarrow R_{\mathrm{cum}} \cdot R(r_2, \beta_2, \gamma_2)
                            = R(r_1, \beta_1, \gamma_1) \cdot R(r_2, \beta_2, \gamma_2).
\]

\textbf{(c)} Compute the joint axis direction, using the \emph{updated}
$R_{\mathrm{cum}}$:
\[
  \unit{a}_2 = R_{\mathrm{cum}} \cdot \unit{a}_2^{\mathrm{local}}.
\]

\paragraph{Iterating.} Continue this three-substep process for each subsequent
joint along the kinematic chain. The output of Step 1 is a list of
$(\vect{p}_i, \unit{a}_i)$ pairs, one per joint, all in the base frame.

\begin{keybox}
\textbf{The key rule.} At each joint, the translation $\vect{t}_i$ is rotated
by the \emph{parent's} cumulative rotation (because $\vect{t}_i$ is expressed
in the parent's frame), while the axis $\unit{a}_i^{\mathrm{local}}$ is
rotated by the \emph{child's} cumulative rotation (because the axis is
expressed in the child's frame, which includes the joint's own rpy).
\end{keybox}

\section{Step 2: Choose Frame 0}

Before computing any DH row, we establish frame 0 --- the base frame of the
DH chain. Frame 0 has two constraints and one free choice:
\begin{itemize}
  \item $\unit{z}_0 = \unit{a}_1$ (required --- frame 0's $z$-axis must lie
        along joint 1's axis).
  \item $\vect{o}_0$ must lie on joint 1's axis. Standard choice:
        $\vect{o}_0 = (0, 0, 0)$ if the URDF's base is already at the origin.
  \item $\unit{x}_0$ is free, as long as it is perpendicular to $\unit{z}_0$.
        Standard choice: project world $+x$ onto the plane perpendicular to
        $\unit{z}_0$, then normalize. If $\unit{z}_0 = (0, 0, 1)$, this gives
        $\unit{x}_0 = (1, 0, 0)$.
\end{itemize}
Complete the right-handed triad with $\unit{y}_0 = \unit{z}_0 \times \unit{x}_0$.

\section{Step 3: Build Frame $i$ from the Common Normal}

Frame $i$ is constructed from the common normal between axis $i$ and axis
$i{+}1$. We describe the process for joint 1 below; the procedure is
identical for each subsequent joint with indices shifted.

\paragraph{Step 3a: Classify the axis pair.} Compute
\[
  \vect{c} = \unit{a}_1 \times \unit{a}_2,
  \qquad \vect{d}_{\mathrm{pp}} = \vect{p}_2 - \vect{p}_1.
\]
\begin{itemize}
  \item If $\norm{\vect{c}} > \epsilon$: axes are skew or intersecting
        (Case~A / Case~B).
  \item If $\norm{\vect{c}} \le \epsilon$: axes are parallel (Case~C).
\end{itemize}
Use $\epsilon \approx 10^{-6}$ in practice.

\paragraph{Step 3b: Find the feet of the common normal.}

\textbf{Cases A and B (non-parallel).} We require the connecting vector
$\vect{f}_2 - \vect{f}_1$ to be perpendicular to both axes. Writing
$\vect{f}_1 = \vect{p}_1 + s\,\unit{a}_1$ and
$\vect{f}_2 = \vect{p}_2 + t\,\unit{a}_2$, the two perpendicularity conditions
\[
  (\vect{f}_2 - \vect{f}_1) \cdot \unit{a}_1 = 0, \qquad
  (\vect{f}_2 - \vect{f}_1) \cdot \unit{a}_2 = 0
\]
give the $2 \times 2$ system:
\[
  \begin{bmatrix} 1 & -\unit{a}_1 \cdot \unit{a}_2 \\
                  \unit{a}_1 \cdot \unit{a}_2 & -1 \end{bmatrix}
  \begin{bmatrix} s \\ t \end{bmatrix}
  = \begin{bmatrix} \vect{d}_{\mathrm{pp}} \cdot \unit{a}_1 \\
                    \vect{d}_{\mathrm{pp}} \cdot \unit{a}_2 \end{bmatrix}.
\]
Solve for $s, t$, then
\[
  \vect{f}_1 = \vect{p}_1 + s\,\unit{a}_1, \qquad
  \vect{f}_2 = \vect{p}_2 + t\,\unit{a}_2.
\]

\textbf{Case C (parallel).} The common normal is non-unique. By convention,
project $\vect{p}_2$ onto axis 1:
\[
  \vect{f}_1 = \vect{p}_1 + (\vect{d}_{\mathrm{pp}} \cdot \unit{a}_1)\,\unit{a}_1,
  \qquad \vect{f}_2 = \vect{p}_2.
\]

\paragraph{Step 3c: Determine $\unit{x}_1$.}

\textbf{Case A (skew, $\norm{\vect{f}_2 - \vect{f}_1} > \epsilon$):}
\[
  \unit{x}_1 = \frac{\vect{f}_2 - \vect{f}_1}{\norm{\vect{f}_2 - \vect{f}_1}}.
\]

\textbf{Case B (intersecting, feet coincide):}
\[
  \unit{x}_1 = \frac{\unit{a}_1 \times \unit{a}_2}{\norm{\unit{a}_1 \times \unit{a}_2}}.
\]
There is still a sign ambiguity; choose the sign so
$\unit{x}_1 \cdot \unit{x}_0 \ge 0$ (this tends to keep $\theta_1$ near zero).

\textbf{Case C (parallel):}
\[
  \unit{x}_1 = \frac{\vect{p}_2 - \vect{f}_1}{\norm{\vect{p}_2 - \vect{f}_1}}.
\]
If $\vect{p}_2 = \vect{f}_1$ (the axes are coincident, not merely parallel),
this formula is degenerate. Fall back to $\unit{x}_0$ projected into the
plane perpendicular to $\unit{a}_2$, then normalized.

\paragraph{Step 3d: Assemble frame 1.}
\[
  \vect{o}_1 = \vect{f}_1, \qquad
  \unit{z}_1 = \unit{a}_2, \qquad
  \unit{x}_1 = \text{as computed above}, \qquad
  \unit{y}_1 = \unit{z}_1 \times \unit{x}_1.
\]
From this point forward, $\unit{z}_1$ refers to the DH frame axis (which
equals $\unit{a}_2$), not the joint 1 axis.

\section{Step 4: Compute the DH Parameters for Joint 1}

With frame 0 (Step 2) and frame 1 (Step 3) both fully defined, apply the
four formulas:
\begin{align}
  a_1      &= (\vect{f}_2 - \vect{f}_1) \cdot \unit{x}_1 \\[4pt]
  \alpha_1 &= \atantwo\!\Bigl(\bigl(\unit{z}_0 \times \unit{z}_1\bigr) \cdot \unit{x}_1,\;\;
                               \unit{z}_0 \cdot \unit{z}_1\Bigr) \\[4pt]
  d_1      &= (\vect{o}_1 - \vect{o}_0) \cdot \unit{z}_0 \\[4pt]
  \theta_1 &= \atantwo\!\Bigl(\bigl(\unit{x}_0 \times \unit{x}_1\bigr) \cdot \unit{z}_0,\;\;
                               \unit{x}_0 \cdot \unit{x}_1\Bigr)
\end{align}
For a revolute joint, $\theta_1$ is a constant offset; the runtime joint
angle is $\theta_1 + q_1$, where $q_1$ is the commanded joint variable. For
a prismatic joint, $d_1$ is the constant offset and the runtime distance is
$d_1 + q_1$.

\subsection*{Validation}

Before trusting the result, check:
\begin{itemize}
  \item $a_1 \ge 0$ in Cases A and B (by construction, if $\unit{x}_1$ was
        chosen as the foot-to-foot direction).
  \item Compose $T_1^0 \, T_2^1 \cdots T_n^{n-1}$ at the zero configuration
        and confirm the resulting end-effector pose matches the URDF forward
        kinematics. Any mismatch signals a sign error somewhere in the chain.
\end{itemize}

\section{Step 5: Iterate to the Next Joint}

To compute joint 2's DH row, repeat Steps 3 and 4:
\begin{itemize}
  \item Use frame 1 (just built) in place of frame 0.
  \item Build frame 2 from the common normal between axis 2 ($\unit{a}_2$)
        and axis 3 ($\unit{a}_3$).
  \item Apply the four DH formulas with the indices shifted by one.
\end{itemize}
Continue until all $n$ joints have been processed.

\begin{warnbox}
\textbf{The last joint.} For joint $n$, there is no axis $n{+}1$. The
$z$-axis of frame $n$ is still constrained to lie along joint $n$'s axis
($\unit{z}_n = \unit{a}_n$), but the origin placement along that axis and
the direction of $\unit{x}_n$ are free. Standard practice is to place
$\vect{o}_n$ at the tool center point and choose $\unit{x}_n$ to match the
tool frame's $x$-axis, both supplied by the URDF's \texttt{tool0} fixed
joint.
\end{warnbox}

\section{Worked Example: Joint 1 of the Motoman GP12}

This example illustrates \textbf{Case A (skew axes)}.

The URDF declares:
\begin{verbatim}
<joint name="joint_1_s" type="revolute">
  <origin rpy="0 0 0" xyz="0 0 0.450"/>
  <axis xyz="0 0 1"/>
</joint>
<joint name="joint_2_l" type="revolute">
  <origin rpy="0 0 0" xyz="0.155 0 0"/>
  <axis xyz="0 1 0"/>
</joint>
\end{verbatim}

\paragraph{Step 1 (axis data).} All \texttt{rpy} values are zero, so
$R_{\mathrm{cum}} = I$ throughout and axis directions pass through unchanged.
Cumulative positions simplify to the sum of \texttt{xyz} translations:
\begin{align*}
  \vect{p}_1 &= (0, 0, 0.450), & \unit{a}_1 &= (0, 0, 1) \\
  \vect{p}_2 &= (0.155, 0, 0.450), & \unit{a}_2 &= (0, 1, 0)
\end{align*}

\paragraph{Step 2 (frame 0).}
\[
  \vect{o}_0 = (0, 0, 0), \quad
  \unit{z}_0 = (0, 0, 1), \quad
  \unit{x}_0 = (1, 0, 0), \quad
  \unit{y}_0 = (0, 1, 0).
\]

\paragraph{Step 3a (classification).}
$\vect{c} = \unit{a}_1 \times \unit{a}_2 = (0,0,1) \times (0,1,0) = (-1, 0, 0)$,
with $\norm{\vect{c}} = 1 > \epsilon$. Case A (skew).

\paragraph{Step 3b (feet of the common normal).}
$\vect{d}_{\mathrm{pp}} = (0.155, 0, 0)$, $\unit{a}_1 \cdot \unit{a}_2 = 0$,
so the system becomes:
\[
  \begin{bmatrix} 1 & 0 \\ 0 & -1 \end{bmatrix}
  \begin{bmatrix} s \\ t \end{bmatrix}
  = \begin{bmatrix} 0 \\ 0 \end{bmatrix}
  \quad\Longrightarrow\quad s = 0,\; t = 0.
\]
\[
  \vect{f}_1 = (0, 0, 0.450), \qquad \vect{f}_2 = (0.155, 0, 0.450).
\]

\paragraph{Step 3c ($\unit{x}_1$).}
$\unit{x}_1 = (0.155, 0, 0) / 0.155 = (1, 0, 0)$.

\paragraph{Step 3d (assemble frame 1).}
\[
  \vect{o}_1 = (0, 0, 0.450), \quad
  \unit{z}_1 = \unit{a}_2 = (0, 1, 0), \quad
  \unit{x}_1 = (1, 0, 0), \quad
  \unit{y}_1 = (0, 0, -1).
\]

\paragraph{Step 4 (DH parameters).}
\begin{align*}
  a_1 &= (\vect{f}_2 - \vect{f}_1) \cdot \unit{x}_1
       = (0.155, 0, 0) \cdot (1, 0, 0) = 0.155 \\[4pt]
  \alpha_1 &= \atantwo\bigl((\unit{z}_0 \times \unit{z}_1) \cdot \unit{x}_1,\;
                              \unit{z}_0 \cdot \unit{z}_1\bigr) \\
           &= \atantwo\bigl((-1, 0, 0) \cdot (1, 0, 0),\; 0\bigr)
            = \atantwo(-1,\, 0) = -\tfrac{\pi}{2} \\[4pt]
  d_1 &= (\vect{o}_1 - \vect{o}_0) \cdot \unit{z}_0
       = (0, 0, 0.450) \cdot (0, 0, 1) = 0.450 \\[4pt]
  \theta_1 &= \atantwo\bigl((\unit{x}_0 \times \unit{x}_1) \cdot \unit{z}_0,\;
                              \unit{x}_0 \cdot \unit{x}_1\bigr)
            = \atantwo(0,\, 1) = 0
\end{align*}

\begin{keybox}
\textbf{Result for joint 1:}
\[
  a_1 = 0.155\,\text{m}, \quad
  \alpha_1 = -\tfrac{\pi}{2}, \quad
  d_1 = 0.450\,\text{m}, \quad
  \theta_1 = 0.
\]
\end{keybox}

\section{Worked Example: Joint 2 of the Motoman GP12}

This example illustrates \textbf{Case C (parallel axes)} and reuses frame 1
as computed in the previous example.

The URDF declares:
\begin{verbatim}
<joint name="joint_2_l" type="revolute">
  <origin rpy="0 0 0" xyz="0.155 0 0"/>
  <axis xyz="0 1 0"/>
</joint>
<joint name="joint_3_u" type="revolute">
  <origin rpy="0 0 0" xyz="0 0 0.614"/>
  <axis xyz="0 -1 0"/>
</joint>
\end{verbatim}

\paragraph{Step 1 (axis data).} Continuing the cumulative sum:
\begin{align*}
  \vect{p}_2 &= (0.155, 0, 0.450), & \unit{a}_2 &= (0, 1, 0) \\
  \vect{p}_3 &= (0.155, 0, 1.064), & \unit{a}_3 &= (0, -1, 0)
\end{align*}

Note that $\unit{a}_3 = -\unit{a}_2$. The URDF deliberately assigns opposite
signs to the shoulder and elbow axis directions --- physically this doesn't
change anything (an axis line is symmetric under sign flip), but it does
affect the DH table. A positive $\unit{a}_3$ would produce $\alpha_2 = 0$
instead of $\alpha_2 = \pi$; the final forward kinematics would still match.

\paragraph{Frame 1 (from the previous example).}
\[
  \vect{o}_1 = (0, 0, 0.450), \quad
  \unit{z}_1 = (0, 1, 0), \quad
  \unit{x}_1 = (1, 0, 0), \quad
  \unit{y}_1 = (0, 0, -1).
\]

\paragraph{Step 3a (classification).}
\[
  \vect{c} = \unit{a}_2 \times \unit{a}_3 = (0,1,0) \times (0,-1,0) = (0,0,0),
  \qquad \norm{\vect{c}} = 0.
\]
Case C (parallel). The axes are antiparallel but that is still ``parallel''
for classification: lines in 3D have no inherent direction, so the common
normal is non-unique regardless of the sign of $\unit{a}_3$.

\paragraph{Step 3b (feet of the common normal, parallel convention).} Set
$\vect{f}_3 = \vect{p}_3$ and project $\vect{p}_3$ onto axis 2:
\[
  \vect{d}_{\mathrm{pp}} = \vect{p}_3 - \vect{p}_2 = (0, 0, 0.614),
  \qquad \vect{d}_{\mathrm{pp}} \cdot \unit{a}_2 = 0.
\]
\[
  \vect{f}_2 = \vect{p}_2 + 0 \cdot \unit{a}_2 = (0.155, 0, 0.450),
  \qquad \vect{f}_3 = (0.155, 0, 1.064).
\]

\paragraph{Step 3c ($\unit{x}_2$, parallel convention).}
\[
  \unit{x}_2 = \frac{\vect{p}_3 - \vect{f}_2}{\norm{\vect{p}_3 - \vect{f}_2}}
             = \frac{(0, 0, 0.614)}{0.614} = (0, 0, 1).
\]

\paragraph{Step 3d (assemble frame 2).}
\[
  \vect{o}_2 = (0.155, 0, 0.450), \quad
  \unit{z}_2 = \unit{a}_3 = (0, -1, 0), \quad
  \unit{x}_2 = (0, 0, 1), \quad
  \unit{y}_2 = \unit{z}_2 \times \unit{x}_2 = (-1, 0, 0).
\]

\paragraph{Step 4 (DH parameters).}
\begin{align*}
  a_2 &= (\vect{f}_3 - \vect{f}_2) \cdot \unit{x}_2
       = (0, 0, 0.614) \cdot (0, 0, 1) = 0.614 \\[4pt]
  \alpha_2 &= \atantwo\bigl((\unit{z}_1 \times \unit{z}_2) \cdot \unit{x}_2,\;
                              \unit{z}_1 \cdot \unit{z}_2\bigr) \\
           &= \atantwo\bigl(\,\vect{0} \cdot (0,0,1),\; -1\,\bigr)
            = \atantwo(0,\, -1) = \pi \\[4pt]
  d_2 &= (\vect{o}_2 - \vect{o}_1) \cdot \unit{z}_1
       = (0.155, 0, 0) \cdot (0, 1, 0) = 0 \\[4pt]
  \theta_2 &= \atantwo\bigl((\unit{x}_1 \times \unit{x}_2) \cdot \unit{z}_1,\;
                              \unit{x}_1 \cdot \unit{x}_2\bigr) \\
           &= \atantwo\bigl((0, -1, 0) \cdot (0, 1, 0),\; 0\bigr)
            = \atantwo(-1,\, 0) = -\tfrac{\pi}{2}
\end{align*}

\begin{keybox}
\textbf{Result for joint 2:}
\[
  a_2 = 0.614\,\text{m}, \quad
  \alpha_2 = \pi, \quad
  d_2 = 0, \quad
  \theta_2 = -\tfrac{\pi}{2}.
\]
\end{keybox}

\paragraph{Interpretation of $\theta_2 = -\pi/2$.} The URDF's zero
configuration places joint 2 ``straight up'' (link 2 oriented along world
$+z$), but the DH frame convention puts $\unit{x}_2$ along the common
normal --- which also points along world $+z$. The $-\pi/2$ offset rotates
$\unit{x}_1$ (world $+x$) to $\unit{x}_2$ (world $+z$) about $\unit{z}_1$
(world $+y$). At runtime, the commanded angle $q_2$ is added to this
offset: $\theta_2^{\mathrm{runtime}} = -\pi/2 + q_2$.

This offset is \emph{not} a bug --- it is an artifact of the URDF and DH
conventions disagreeing about what the zero pose looks like. Published DH
tables for the GP12 typically bake this offset into the joint angle
definition, so the ``canonical'' and ``URDF'' zero poses differ by exactly
this quarter turn.

\section{Sign-Convention Pitfalls}

The most common errors in DH derivations come from sign ambiguities. Three
rules prevent most of them:

\begin{enumerate}
  \item \textbf{The four DH parameters are not independent scalars --- they
    are the components of a rigid transform.} Flipping the sign of
    $\unit{x}_i$ flips both $a_i$ and $\alpha_i$, and shifts $\theta_i$ by
    $\pi$. You cannot flip one in isolation.
  \item \textbf{Commit to a convention for $\unit{x}_i$ before computing.}
    The ``common normal points from axis $i$ to axis $i{+}1$'' convention
    (used above) gives $a_i \ge 0$ and tends to produce clean tables.
  \item \textbf{Validate by composing the transforms.} At the zero
    configuration, the product $T_1^0 \cdots T_n^{n-1}$ must match the
    end-effector pose computed directly from the URDF. If it doesn't, there
    is a sign error, and it is almost always in $\alpha_i$ or $\theta_i$.
\end{enumerate}

\end{document}
